Dirichlet 前缀和（也称“倍数前缀和”）是处理数论中因子/倍数关系的重要工具。对于函数 $f(n)$，定义其 Dirichlet 前缀和为：
\[
g(n) = \sum_{d\mid n} f(d),
\]
即 $g$ 是 $f$ 在所有因子上的求和。类似地，Dirichlet 后缀和为：
\[
h(n) = \sum_{n\mid d} f(d),
\]
即 $h$ 是 $f$ 在所有倍数上的求和。其逆运算（差分）可通过 Möbius 反演实现：
\[
f(n) = \sum_{d\mid n} \mu(d) \, g\!\left(\frac{n}{d}\right), \quad
f(n) = \sum_{n\mid d} \mu\!\left(\frac{d}{n}\right) h(d).
\]

直接对每个 $n$ 枚举因子或倍数的时间复杂度为 $O(n \log n)$（调和级数）。  
通过线性筛预处理所有质数，可以 $O(n \log \log n)$ 完成 Dirichlet 前缀和/后缀和，其思想与高维前缀和类似：将每个数看作一个高维向量（质因子的指数），倍数关系就是每一维的指数均不小于原数。因此，对每个质数 $p$，按指数从低到高做一次“前缀和”即可。